\documentclass[11pt]{article}

\usepackage[margin=1in]{geometry}
\usepackage{amsmath,amssymb,amsthm,mathtools}
\usepackage{booktabs,longtable,array,tabularx}
\usepackage{enumitem}
\usepackage{microtype}
\usepackage{hyperref}
\usepackage{xcolor}

\hypersetup{
  colorlinks=true,
  linkcolor=blue,
  citecolor=blue,
  urlcolor=blue
}

\newtheorem{theorem}{Theorem}[section]
\newtheorem{proposition}[theorem]{Proposition}
\newtheorem{lemma}[theorem]{Lemma}
\newtheorem{corollary}[theorem]{Corollary}
\theoremstyle{remark}
\newtheorem{remark}[theorem]{Remark}

\newcommand{\N}{\mathbb{N}}
\newcommand{\Nzero}{\mathbb{N}_0}
\newcommand{\code}[1]{\texttt{\detokenize{#1}}}
\newcommand{\cpath}[1]{\path{#1}}
\newcommand{\ARCKernel}{\mathcal{K}}
\newcommand{\Change}{\mathcal{D}}

\title{\textbf{Cobham-Free Full General Automatic Rigidity\\
for Collatz-Invariant Sequences}\\[0.4em]
\large Removing the Final External Cobham Dependency in the ARC Program}

\author{Muteki\\
Independent Researcher, Japan}

\date{\textbf{First released: September 23, 2026}}

\begin{document}

\maketitle

\begin{abstract}
We study finite-kernel automatic colorings of the nonnegative integers that are
invariant under the shortcut Collatz map
\[
T(n)=
\begin{cases}
n/2, & n\equiv 0\pmod 2,\\[2mm]
(3n+1)/2, & n\equiv 1\pmod 2.
\end{cases}
\]
Earlier stages of the ARC (Automatic Rigidity for Collatz) program established
positive-domain rigidity for several individual bases, then for every odd base,
and finally for all bases under an explicit Cobham-theorem handoff localized to
the base-$2$ / pure-power-of-two branch.  The present paper removes that final
external dependency.

The key new step is a Cobham-free reconstruction of ARC2.  For a base-$2$
finite-kernel automatic invariant coloring $a$, we pass to the odd subsequence
$b(n)=a(2n+1)$.  Existing ARC2 kernel arithmetic makes $b$ base-$3$
finite-kernel automatic.  Instead of applying Cobham's theorem to obtain
ultimate periodicity, local-density machinery developed later in the ARC
program shows that the adjacent-change set of $b$ is dyadically nowhere thick.
The General Odd ARC topological-pumping theorem then forces that change set to
be finite.  Hence $b$ is eventually constant, and the recurrence
$b(4n+2)=b(n)$ propagates the tail value to every index.  Doubling invariance
then recovers constancy of $a$ on all positive integers.

Combining this new ARC2 theorem with odd-base rigidity and base-factor descent
yields a single theorem for every base $B\ge 2$, with no explicit
\texttt{ARCCobhamPrinciple} hypothesis and no \texttt{[Finite alpha]}
typeclass assumption.  The Lean 4 formalization was independently re-derived
in a final audit; the final dependency report contains only
\texttt{propext}, \texttt{Classical.choice}, and \texttt{Quot.sound}, with no
\texttt{sorryAx}.  The integrated project completed \texttt{lake build}
successfully with 8906 jobs.
\end{abstract}

\section{Introduction}

Automatic sequences provide a natural finite-state lens through which to study
arithmetic dynamics.  The ARC program asks a deliberately restricted question:
if a finite-state coloring is invariant along every edge of the shortcut
Collatz graph, how much freedom can remain?

Let
\[
T(2n)=n,\qquad T(2n+1)=3n+2,
\]
which is equivalent on odd inputs to $(3m+1)/2$.  A coloring
$a:\Nzero\to\alpha$ is called \emph{shortcut-Collatz invariant} if
\[
a(T(n))=a(n)\qquad(n\in\Nzero).
\]
The ARC objective is not to prove the Collatz conjecture.  Rather, it is to
show that a broad class of finite-state invariants is too rigid to distinguish
positive integers.

The earlier development proceeded through increasingly general cases.
ARC2 and ARC3 used kernel transfer together with an explicit external form of
Cobham's theorem.  ARC5 and ARC7 introduced local synchronization, dyadic
thinness, pumping, and finite-change arguments.  ARC15 stress-tested the method
at an odd composite base.  These ideas were then generalized to every odd
base.  Base-factor descent subsequently reduced every mixed even base
$B=2^s m$ with odd $m>1$ to its odd part.

At that point, the only remaining explicit Cobham dependency in the integrated
all-base theorem was the pure power-of-two branch $B=2^s$, because factor
descent sends this branch to base $2$, and the old ARC2 endpoint still passed
through Cobham.

The purpose of the present paper is to close precisely that gap.

\begin{theorem}[Cobham-free Full General ARC]\label{thm:main}
Let $B\ge2$ and let $a:\Nzero\to\alpha$ have finite base-$B$ kernel.
Assume
\[
a(T(n))=a(n)\qquad(n\in\Nzero).
\]
Then there exists $c\in\alpha$ such that
\[
a(n)=c\qquad(n>0).
\]
\end{theorem}

The value at zero is not determined; this is sharp.

The proof has three layers:
\begin{enumerate}[label=(\roman*)]
\item reconstruct ARC2 without Cobham;
\item use that reconstruction to close every pure power-of-two base;
\item combine the pure-power branch with the already Cobham-free odd and mixed
      even branches.
\end{enumerate}

\section{Finite-kernel automaticity and shortcut invariance}

For a base $b\ge2$, define the $b$-kernel of a sequence $a:\Nzero\to\alpha$ by
\[
\ARCKernel_b(a)
=
\left\{
n\longmapsto a(b^e n+r)
:
e\ge0,\ 0\le r<b^e
\right\}.
\]
We use the finite-kernel condition
\[
\ARCKernel_b(a)\ \text{is finite}
\]
as the formal notion of base-$b$ automaticity.  This is the interface used by
the principal Lean theorems.

Shortcut invariance immediately yields the branch identities
\begin{equation}\label{eq:branches}
a(2n)=a(n),
\qquad
a(2n+1)=a(3n+2).
\end{equation}
The first identity is especially important: for every $k\ge0$,
\begin{equation}\label{eq:doubling}
a(2^k n)=a(n).
\end{equation}

Equation \eqref{eq:doubling} is the source of both factor descent and the final
propagation of eventual constancy to all positive indices.

\section{Where the old proof still used Cobham}

Suppose the base is even.  Write
\[
B=2^s m,
\qquad
s>0,
\qquad
m\ \text{odd}.
\]
The ARC base-factor descent theorem shows that, under doubling invariance,
\begin{equation}\label{eq:descent}
(2^s m)\text{-finite-kernel automatic}
\quad\Longrightarrow\quad
m\text{-finite-kernel automatic}.
\end{equation}

If $m>1$, then $m\ge3$ is odd.  The completed General Odd ARC theorem already
gives positive-domain constancy directly, so this branch is Cobham-free.

The only unresolved case after localization was
\[
m=1,
\qquad\text{i.e.}\qquad
B=2^s.
\]
Applying \eqref{eq:descent} repeatedly reduces this branch to base $2$.
Thus the all-base proof becomes Cobham-free as soon as ARC2 itself is
reconstructed without Cobham.

This observation is conceptually useful: the problem is not ``remove Cobham
everywhere.''  By the time of the present work, the dependency had already been
localized to one sharply identified endpoint.

\section{Cobham-free reconstruction of ARC2}

Let $a:\Nzero\to\alpha$ be base-$2$ finite-kernel automatic and
shortcut-Collatz invariant.  Define the odd subsequence
\begin{equation}\label{eq:bdef}
b(n):=a(2n+1).
\end{equation}

The new proof does not try to prove that $b$ is shortcut-Collatz invariant.
Instead, it extracts exactly the three properties needed for rigidity:
base-$3$ finite-kernel automaticity, dyadic nowhere-thickness of the adjacent
change set, and the self-similarity relation $b(4n+2)=b(n)$.

\subsection{The odd subsequence is base 3 automatic}

The original ARC2 arithmetic already contains the essential kernel transfer.
Using \eqref{eq:branches}, every relevant base-$3$ decimation of the odd
subsequence can be represented inside the finite base-$2$ kernel generated by
$a$.  In particular,
\[
\ARCKernel_3(b)
\]
is finite.

Formally, the old lower-level theorem
\code{arc_three_kernel_odd_finite}
is reused directly.  The new wrapper
\[
\code{arc2OddSubseq_threeAutomatic_of_twoAutomatic}
\]
works purely at finite-kernel level and avoids the unnecessary
\code{[Finite alpha]} hypothesis carried by an older high-level interface.

\begin{proposition}\label{prop:auto3}
If $a$ has finite base-$2$ kernel and satisfies shortcut invariance, then
$b(n)=a(2n+1)$ has finite base-$3$ kernel.
\end{proposition}

This is an important point in the formal architecture: no new arithmetic
kernel-transfer theorem was needed.  What changed was the way the existing
transfer theorem was connected to the later topological machinery.

\subsection{The adjacent-change set is dyadically nowhere thick}

Define
\[
\Change_b
=
\{n\in\Nzero : b(n)\ne b(n+1)\}.
\]
Since
\[
b(n)=a(2n+1),
\qquad
b(n+1)=a(2n+3),
\]
the equality $b(n)=b(n+1)$ is equivalent to
\begin{equation}\label{eq:gap2}
a(2n+1)=a(2n+3).
\end{equation}

The decisive new observation is that the local-density theorem for
shortcut-invariant colorings can be applied with gap $H=2$.

Informally, fix any dyadic parent cylinder
\[
n\equiv r\pmod{2^m}.
\]
Under the affine map $x=2n+1$, this becomes the odd dyadic cylinder
\[
x\equiv 2r+1\pmod{2^{m+1}}.
\]
The local-density theorem supplies a deeper subcylinder on which
\[
a(x)=a(x+2).
\]
Pulling this subcylinder back through $x=2n+1$ gives a deeper dyadic child in
the original $n$-variable on which
\[
b(n)=b(n+1).
\]
Therefore that child contains no element of $\Change_b$.

\begin{proposition}[Dyadic thinness]\label{prop:thin}
For every dyadic parent cylinder there is a deeper dyadic child disjoint from
$\Change_b$.  Equivalently, $\Change_b$ is dyadically nowhere thick.
\end{proposition}

In Lean this is established through
\[
\code{arc2OddSubseq_local_gap_two_child},
\qquad
\code{arc2OddSubseq_changeSet_nowhereThick},
\]
and packaged as
\[
\code{arc2CobhamRemoval_stage1_thinness}.
\]

This is the step that replaces the role previously played by Cobham's theorem.
The old proof obtained global regularity through ultimate periodicity.
The new proof obtains a different form of global control: automaticity plus
dyadic topological thinness forces only finitely many adjacent changes.

\subsection{Finite change set and eventual constancy}

The General Odd ARC development proves the required topological-pumping
principle for odd automatic bases.  Applied at base $3$, it says in the present
setting that a base-$3$ finite-kernel automatic adjacent-change set that is
dyadically nowhere thick must be finite.

Combining Proposition~\ref{prop:auto3} and Proposition~\ref{prop:thin} gives
\[
\Change_b\ \text{finite}.
\]
Hence there exists $N$ such that
\[
b(n)=b(n+1)\qquad(n\ge N),
\]
and therefore $b$ is eventually constant.

The relevant reusable General Odd ARC ingredients are the change-set
automaticity, nowhere-thick-to-finite support theorem, and finite-change-to-tail
constancy theorem.

\subsection{Tail constancy propagates to every index}

Shortcut invariance also gives the exact recurrence
\begin{equation}\label{eq:bstep}
b(4n+2)=b(n).
\end{equation}
Indeed,
\[
b(4n+2)
=
a(8n+5),
\]
and repeated use of the shortcut branch relations returns the value to
$a(2n+1)=b(n)$.

Now iterate the affine map
\[
\Phi(n)=4n+2.
\]
For every fixed $n$, the values
\[
n,\ \Phi(n),\ \Phi^2(n),\ldots
\]
grow without bound, while \eqref{eq:bstep} preserves the color.  Thus, once an
iterate lies beyond the eventual-constancy threshold,
\[
b(n)=c.
\]
So $b$ is constant on all of $\Nzero$.

Finally, every positive integer can be written as
\[
n=2^k q
\]
with $q$ odd.  By \eqref{eq:doubling},
\[
a(n)=a(q),
\]
and $a(q)$ is a value of the odd subsequence $b$.  Hence $a$ is constant on all
positive integers.

\begin{theorem}[Cobham-free ARC2]\label{thm:arc2}
Let $a:\Nzero\to\alpha$ have finite base-$2$ kernel and satisfy shortcut
invariance.  Then there exists $c\in\alpha$ such that
\[
a(n)=c\qquad(n>0).
\]
\end{theorem}

The principal Lean endpoint is
\[
\code{arc2_main_kernel_form_no_cobham},
\]
with stage wrapper
\[
\code{arc2CobhamRemoval_stage2_final}.
\]

\section{Closing every base}

With Theorem~\ref{thm:arc2} available, the last missing branch of the general
argument becomes elementary.

\begin{proposition}[Pure powers of two]\label{prop:pure}
Let
\[
B=2^s,\qquad s>0.
\]
If $a$ is base-$B$ finite-kernel automatic and shortcut-Collatz invariant,
then $a$ is constant on the positive integers.
\end{proposition}

\begin{proof}
Write $s=t+1$.  Base-factor descent gives base-$2$ finite-kernel automaticity.
Apply Theorem~\ref{thm:arc2}.
\end{proof}

In Lean, this is separated into
\[
\code{arc_powTwoSucc_final_no_cobham}
\]
and
\[
\code{arc_powTwoBase_final_no_cobham}.
\]

\begin{proposition}[Mixed even bases]\label{prop:mixed}
Let
\[
B=2^s m,
\qquad
s>0,
\qquad
m>1\ \text{odd}.
\]
If $a$ is base-$B$ finite-kernel automatic and shortcut-Collatz invariant,
then $a$ is constant on the positive integers.
\end{proposition}

\begin{proof}
Factor descent gives finite base-$m$ kernel.  Since $m\ge3$ is odd, apply the
General Odd ARC theorem.
\end{proof}

\begin{proof}[Proof of Theorem~\ref{thm:main}]
If $B$ is odd, then $B\ge3$ and General Odd ARC applies directly.

If $B$ is even, factor
\[
B=2^s m
\]
with $s>0$ and $m$ odd.  If $m=1$, use
Proposition~\ref{prop:pure}.  If $m>1$, use
Proposition~\ref{prop:mixed}.  These cases exhaust all $B\ge2$.
\end{proof}

The packaged Lean endpoint is
\[
\code{arc_full_general_no_cobham}.
\]

\section{Lean 4 formalization}

The Cobham-removal layer was developed conservatively: each stage was checked
standalone, integrated into the root import graph, and followed by a project
build before proceeding.

\begin{center}
\footnotesize
\renewcommand{\arraystretch}{1.16}
\begin{tabularx}{\textwidth}{@{}p{0.27\textwidth}X p{0.24\textwidth}@{}}
\toprule
\textbf{File} & \textbf{Role} & \textbf{Principal endpoint}\\
\midrule

\texttt{ARC2CobhamRemoval}\newline\texttt{Stage0.lean}
&
Odd-subsequence packaging, recurrence extraction, eventual-to-global
propagation, and recovery of positive-domain constancy.
&
\cpath{arc2OddSubseq_rigidity_of_threeAutomatic_and_thin}
\\[0.7em]

\texttt{ARC2CobhamRemoval}\newline\texttt{Stage1\_Thinness\_v3.lean}
&
Derives dyadic nowhere-thickness of the odd-subsequence adjacent-change set
from shortcut invariance and local density.
&
\cpath{arc2CobhamRemoval_stage1_thinness}
\\[0.7em]

\texttt{ARC2CobhamRemoval}\newline\texttt{Stage2\_Final\_v2.lean}
&
Reconnects directly to the lower ARC2 finite-kernel transfer theorem and
closes ARC2 without Cobham or \code{[Finite alpha]}.
&
\cpath{arc2CobhamRemoval_stage2_final}
\\[0.7em]

\texttt{ARC2CobhamRemoval}\newline\texttt{Stage3\_FullGeneral.lean}
&
Replaces the pure power-of-two branch, reassembles all even bases, and closes
every base.
&
\cpath{arc_full_general_no_cobham}
\\[0.7em]

\texttt{ARCFullGeneralNoCobham}\newline\texttt{FinalAudit.lean}
&
Independently reconstructs the all-base theorem from lower-level components
and runs representative smoke tests.
&
\cpath{arcFinalAudit_full_general_no_cobham}
\\

\bottomrule
\end{tabularx}
\end{center}

\subsection{Final audit}

The final audit deliberately does not prove the all-base result by merely
calling the packaged Stage-3 endpoint.  It reconstructs the three structural
branches:
\[
\text{pure }2^s,
\qquad
\text{mixed even},
\qquad
\text{odd},
\]
from lower-level theorems.

Representative concrete instances were checked at bases
\[
2,3,4,5,6,7,10,15,64,100.
\]

For the principal final theorems,
\code{\#print axioms} reports
\[
\code{propext},
\qquad
\code{Classical.choice},
\qquad
\code{Quot.sound},
\]
and does not report \code{sorryAx}.  No explicit
\code{ARCCobhamPrinciple} parameter occurs in the final theorem, and no
\code{[Finite alpha]} typeclass assumption is required.

After the final audit import was added to the root project, the command
\[
\code{lake build}
\]
completed successfully with \textbf{8906 jobs}.

\begin{remark}[Scope of the formal audit]
The Lean dependency report verifies the proof as encoded relative to the
imported definitions and Mathlib environment.  It is not, by itself, an
independent peer review of the mathematical modeling choices or a literature
novelty certificate.
\end{remark}

\section{Relation to the earlier ARC papers}

The present release is intentionally additive rather than substitutive.

The earlier ARC papers are retained because they document the route by which
the final architecture emerged:
\begin{itemize}
\item ARC2 contains the original base-$2$ kernel arithmetic and the first
      rigidity route.
\item ARC3 develops a complementary kernel-transfer case.
\item ARC5 introduces local synchronization, dyadic thinness, and a
      finite-change strategy.
\item ARC7 develops doubled-loop pumping and dyadic filling.
\item ARC15 stress-tests the mechanism at an odd composite base.
\item General Odd ARC extracts the common odd-base structure.
\item The all-even and first Full General integrations identify exactly where
      the external Cobham handoff remained.
\end{itemize}

The new proof feeds machinery developed later in the program back into ARC2.
That feedback loop is the conceptual content of the present paper.

Thus this work should not be described as correcting an error in the earlier
papers.  The earlier theorems remain valid under their stated proof
architectures.  The present result removes an external hypothesis from the
integrated formal route and therefore strengthens the final theorem.

\section{Sharpness at zero}

The restriction to positive integers cannot be removed.

\begin{proposition}\label{prop:zero}
Let $x,y\in\alpha$ be arbitrary and define
\[
a(0)=x,
\qquad
a(n)=y\quad(n>0).
\]
Then $a$ is finite-kernel automatic in every base and shortcut-Collatz
invariant.
\end{proposition}

Consequently, the value at zero is genuinely free.  Theorem~\ref{thm:main}
is sharp in its domain of constancy.

For Boolean colorings, Theorem~\ref{thm:main} may equivalently be read as a
rigidity statement for automatic invariant subsets of $\Nzero$: on the
positive integers, such a set is either empty or all of $\N_{>0}$, while
membership of zero remains independent.

\section{What the theorem does not say}

It is important to separate the ARC theorem from the Collatz conjecture itself.

Theorem~\ref{thm:main} says that a finite-kernel automatic coloring cannot
provide a nontrivial positive-domain invariant of the shortcut Collatz map.
It does not assert that every Collatz orbit reaches $1$, nor does it rule out
all conceivable invariants.  In particular, arbitrary nonautomatic or
infinite-state colorings lie outside the theorem's hypotheses.

Similarly, removing Cobham from the present proof does not diminish the
importance of Cobham's theorem in automatic-sequence theory.  The result is
specific to the additional arithmetic and dynamical structure supplied by
shortcut-Collatz invariance.

\section{Release note}

\textbf{First released: September 23, 2026.}

Earlier ARC papers are not superseded or withdrawn.  They are retained as
historical stages of the development.  The present paper is added as a new
release documenting the Cobham-free closure of the Full General ARC theorem.

\section*{Acknowledgment of AI assistance}

OpenAI ChatGPT (GPT-5.6 Sol) was used as a research and drafting aid for
mathematical exploration, theorem organization, Lean 4 code development and
debugging, formal dependency auditing, manuscript preparation, and repository
organization.  The AI system is not an author or bibliographic source.
Responsibility for the released mathematical claims, formalization, references,
and wording remains with the human author.

\begin{thebibliography}{9}

\bibitem{AlloucheShallit}
J.-P. Allouche and J. Shallit,
\emph{Automatic Sequences: Theory, Applications, Generalizations},
Cambridge University Press, 2003.

\bibitem{Cobham1969}
A. Cobham,
On the base-dependence of sets of numbers recognizable by finite automata,
\emph{Mathematical Systems Theory} \textbf{3} (1969), 186--192.

\bibitem{Eilenberg}
S. Eilenberg,
\emph{Automata, Languages, and Machines, Vol. A},
Academic Press, 1974.

\end{thebibliography}

\end{document}